\subsection{PROPERTIES OF INFINITE CARDINALS}

THEOREM 2. For every infinite cardinal \(\mathfrak{a}\) we have \(\mathfrak{a}^2 = \mathfrak{a}\) .

We shall use the following two lemmas :

\textbf{LEMMA 1. Every infinite set E contains a set equipotent to N.}

There exists a well-ordering relation on E (§ 2, no. 3, Theorem 1), which we shall denote by \(x \leqslant y\) . The hypothesis implies that the well-ordered set E cannot be isomorphic to a segment of N distinct from N, for such a segment is of the form {[}0, n{]} (§ 2, no. 1, Proposition 1) and is therefore finite. It follows that N is isomorphic to a segment of E (§ 2, no. 3, Theorem 3), whence the result.

\textbf{LEMMA 2. The set N × N is equipotent to N.}

Since \(N \times N\) contains the set \(\{0\} \times N\) , which is equipotent to N, we have \(\text{Card}(N) \leqslant \text{Card}(N \times N)\) . To complete the proof it is enough to define an injection f of \(N \times N\) into N. For this purpose we note that there exists an injection \(\varphi\) of \(\mathbf{N}\) into the set of mappings of \(\mathbf{N}\) into \(\mathrm{I} = \{0,1\}\) , obtained as follows: if \(r\) is the least integer such that \(n > 2^r\) , and if \(\sum_{k=0}^{r-1}\varepsilon_k2^{r-k-1}\) is the dyadic expansion of \(n\) ( \(\S 5\) , no. 7), \(\varphi(n)\) is defined to be the sequence \((u_m)_{m\in\mathbf{N}}\) such that \(u_m = \varepsilon_{r-m-1}\) for \(m < r\) and \(u_m = 0\) for \(m \geqslant r\) . Proposition 8 of \(\S 5\) , no. 7 shows that \(\varphi\) is injective. For each pair \((n, n') \in \mathbf{N} \times \mathbf{N}\) we define \(f(n, n')\) as follows: if \(\varphi(n) = (u_m)\) and \(\varphi(n') = (v_m)\) , let \(f(n, n')\) be the integer \(s\) such that \(\varphi(s) = (w_m)\) , where \(w_{2m} = u_m\) and \(w_{2m+1} = v_m\) for all \(m \in \mathbf{N}\) . It is clear that the relation \(f(n, n') = f(n_1, n_1')\) implies \(\varphi(n) = \varphi(n_1)\) and \(\varphi(n') = \varphi(n_1')\) ; hence \((n, n') = (n_1, n_1')\) , and therefore \(f\) is injective.

We come now to the proof of Theorem 2. Let E be a set such that \(\operatorname{Card}(\mathbf{E}) = \mathfrak{a}\) . Let D be a subset of E equipotent to N (Lemma 1). Then there exists a bijection \(\psi_0\) of D onto D × D (Lemma 2). Let \(\mathfrak{M}\) be the set of pairs (X, \(\psi\) ), where X is a subset of E containing D and \(\psi\) is a bijection of X onto X × X which extends \(\psi_0\) . Order the set \(\mathfrak{M}\) by means of the relation

\textbf{``X ⊂ X' and ψ' is an extension of ψ''}

between (X, \(\psi\) ) and (X', \(\psi'\) ). Then it is immediately seen that \(\mathfrak{M}\) is inductive (cf.~§2, no. 4, Example 2). Hence, by Zorn's Lemma (§2, no. 4, Theorem 2), \(\mathfrak{M}\) has a maximal element (F, \(f\) ). We shall show that Card (F) = a, which will suffice to prove the theorem. If Card (F) is not equal to a, let Card (F) = b \textless{} a. Then b = b² and b is infinite, and b ≤ 2b ≤ 3b ≤ b² = b (§ 3, no. 6, Proposition 14); hence 2b = b and 3b = b. From the hypothesis b \textless{} a it follows that Card (E --- F) \textgreater{} b, for otherwise we should have Card (E) ≤ 2b = b, and we have assumed that b \textless{} Card (E). Hence there is a subset Y ⊂ E --- F equipotent to F. Let Z = F ∪ Y; we shall show that there exists a bijection g of Z onto Z × Z which extends f.~We have

\[
\mathbf {Z} \times \mathbf {Z} = (\mathbf {F} \times \mathbf {F}) \cup (\mathbf {F} \times \mathbf {Y}) \cup (\mathbf {Y} \times \mathbf {F}) \cup (\mathbf {Y} \times \mathbf {Y}),
\]

and the four products on the right-hand side are mutually disjoint. Since F and Y are equipotent, we have

\[
\operatorname{Card} (\mathbf {F} \times \mathbf {Y}) = \operatorname{Card} (\mathbf {Y} \times \mathbf {F}) = \operatorname{Card} (\mathbf {Y} \times \mathbf {Y}) = \mathfrak {b} ^ {2} = \mathfrak {b},
\]

whence

\[
\operatorname{Card} ((\mathbf {F} \times \mathbf {Y}) \cup (\mathbf {Y} \times \mathbf {F}) \cup (\mathbf {Y} \times \mathbf {Y})) = 3 \mathfrak {b} = \mathfrak {b}.
\]

There is therefore a bijection \(f_{1}\) of Y onto the set

\[
(\mathbf {F} \times \mathbf {Y}) \cup (\mathbf {Y} \times \mathbf {F}) \cup (\mathbf {Y} \times \mathbf {Y});
\]

the mapping g of Z into \(Z \times Z\) , which is equal to f on F and to \(f_{1}\) on Y, is therefore a bijection which extends f, contrary to the definition of f.~Hence \(\text{Card}(F) = a\) , and the proof is complete.

COROLLARY 1. If \(\mathfrak{a}\) is an infinite cardinal, then \(\mathfrak{a}^n = \mathfrak{a}\) for every integer \(n \geqslant 1\) . By induction on \(n\) .

COROLLARY 2. The product of a finite family \((\mathfrak{a}_i)_{i\in \mathbf{I}}\) of non-zero cardinals, of which the greatest is an infinite cardinal \(\mathfrak{a}\) , is equal to \(\mathfrak{a}\) .

Let \(\mathfrak{b}\) denote the product and let \(n\) be the number of elements in I. Then \(\mathfrak{b} \leqslant \mathfrak{a}^n = \mathfrak{a}\) (§3, no. 6, Proposition 14). On the other hand, since \(\mathfrak{a}_i \geqslant 1\) for all \(i \in \mathbf{I}\) , we have \(\mathfrak{b} \geqslant \mathfrak{a}\) (§3, no. 6, Proposition 14).

COROLLARY 3. Let \(\mathfrak{a}\) be an infinite cardinal and let \((\mathfrak{a}_{\iota})_{\iota \in I}\) be a family of cardinals \(\leqslant \mathfrak{a}\) whose index set \(I\) has a cardinal \(\leqslant \mathfrak{a}\) . Then \(\sum_{\iota \in I} \mathfrak{a}_{\iota} \leqslant \mathfrak{a}\) ; and if \(\mathfrak{a}_{\iota} = \mathfrak{a}\) for at least one index \(\iota \in I\) , then \(\sum_{\iota \in I} \mathfrak{a}_{\iota} = \mathfrak{a}\) .

Let \(\mathfrak{b}\) be the cardinal of I; then we have \(\sum_{i\in I}a_i\leqslant a\mathfrak{b}\leqslant a^2 = a\) (§3, no. 6, Proposition 14), and \(\sum_{i\in I}a_i\geqslant a_x\) for all \(x\in I\) .

COROLLARY 4. If a and b are two non-zero cardinals, one of which is infinite, we have \(a\mathfrak{b} = a + \mathfrak{b} = \sup(a, \mathfrak{b})\) .

This follows directly from Corollaries 2 and 3.

